% titel: Traegheitskraefte und Coriolis-Effekt an Fluessen
% skript: Kap. Rotierende Bezugssysteme, Corioliskraft

Trägheitskräfte (Scheinkräfte) treten nur in beschleunigten oder rotierenden
Bezugssystemen auf. Wichtigster Fall ist die Corioliskraft in erdgebundenen
Systemen; sie erklärt die seitliche Ablenkung grosser Flüsse und die
asymmetrische Uferausbildung.

\begin{werkzeugkasten}[Werkzeuge]
\textbf{Trägheitskraft} (Scheinkraft) $F_\mathrm{tr} = -m\,\vec a_\mathrm{Tr}$
mit der Beschleunigung des Bezugssystems.

\textbf{Corioliskraft} in einem mit $\vec\Omega$ rotierenden System bei
Geschwindigkeit $\vec v$:
\begin{equation}
  \vec F_C = -2m\,\vec\Omega\times\vec v
  \;\Rightarrow\;
  a_C = 2\,\Omega\,v\,\sin\alpha\;,
\end{equation}
$\alpha$ der Winkel zwischen $\vec\Omega$ (Nordrichtung) und $\vec v$. Für eine
horizontale Fliesströmung zählt die Komponente von $\vec\Omega$ senkrecht dazu,
also $\Omega\,\sin\vartheta$ ($\vartheta$ geographische Breite).

\textbf{Bankneigung:} Die durch die Querkraft entstehende Querbeschleunigung
$g$ überlagert sich; der Wasserspiegel stellt sich schief zur
Horizontalen ($\tan\beta = a_C/g$), Höhe über Flussbreite $B$:
\begin{equation}
  \Delta h = B\,\frac{a_C}{g}\;.
\end{equation}

\textbf{Konstanten:} Erdrotation $\Omega\approx\qty{7.27e-5}{rad/s}$,
$g=\qty{9.81}{m/s^2}$.
\end{werkzeugkasten}

\begin{werkzeugkasten}[Vorgehen in allgemeiner Form]
(a) Die drei Unterscheidungsmerkmale von Trägheits- und echten Kräften
nennen. (b) $a_C$ mit der Projektion von $\Omega$ auf die vom Fluss gesehene
Richtung berechnen. (c) über $\Delta h=B\cdot a_C/g$ die Uferhöhe bestimmen und
die Ablenkrichtung begründen.
\end{werkzeugkasten}

\begin{achtung}[Typische Fallen]
\begin{itemize}
  \item $\vec\Omega$ liegt in Richtung der Erdachse (nach Norden); für eine
       horizontale Strömung den $\sin\vartheta$-Faktor der Breite nicht
       vergessen.
  \item Auf der Nordhalbkugel erfolgt die Ablenkung immer nach \emph{rechts}
       der Bewegungsrichtung.
\end{itemize}
\end{achtung}
